\chapter{Instalación y Configuración}
\label{ch:setup}
% ==============================================================

\section{Requisitos del sistema}

\begin{table}[H]
\centering
\caption{Requisitos hardware y software}
\begin{tabular}{L{4cm} L{8cm}}
\toprule
\textbf{Componente} & \textbf{Requisito mínimo} \\
\midrule
SO & Linux (Ubuntu 22.04+) o Windows 10+ \\
Python & 3.11+ \\
RAM & 8 GB (16 GB recomendado para SAM) \\
GPU & CUDA opcional pero recomendada para SAM \\
Almacenamiento & 50 GB para datos (6 cámaras × 24h/día) \\
Node.js & 18+ (para frontend) \\
\bottomrule
\end{tabular}
\end{table}

\section{Instalación del entorno Python}

\begin{codebox}[]
\begin{minted}{bash}
conda create -n cv-lit python=3.11
conda activate cv-lit

# Visión + álgebra
pip install opencv-python numpy scipy pandas

# Backend API
pip install fastapi uvicorn python-multipart

# Segmentación SAM
pip install torch torchvision
pip install git+https://github.com/facebookresearch/segment-anything.git

# Geo
pip install gdal pyproj shapely
\end{minted}
\end{codebox}

\section{Instalación del frontend}

\begin{codebox}[]
\begin{minted}{bash}
cd frontend
npm install
npm run dev       # modo desarrollo
npm run build     # build de producción (dist/)
\end{minted}
\end{codebox}

\section{Configuración de la API Obscape}

Las credenciales de acceso a la API se configuran en \texttt{acces\_api/obscape\_api.py}:

\begin{table}[H]
\centering
\caption{Parámetros de conexión Obscape}
\begin{tabular}{L{4cm} L{8cm}}
\toprule
\textbf{Parámetro} & \textbf{Valor} \\
\midrule
Base URL   & \texttt{https://obscape.com/portal/api/v3/api} \\
Username   & \texttt{fuster} \\
Portal     & \texttt{https://obscape.com/portal} \\
\bottomrule
\end{tabular}
\end{table}

% ==============================================================
\appendix
% ==============================================================

